% Vorlage fuer Abschlussarbeiten am Fachbereich Mathematik der TU Darmstadt.
% Geeignet fuer Bachelorarbeiten, Masterarbeiten und Diplomarbeiten.

%\documentclass[12pt,a4paper,twoside]{report}
\documentclass[12pt,a4paper,twoside]{article}
% Benutzt man die Option draft, so kann man die Umbrueche ueberpruefen.

% Hier werden alle benoetigten Pakete und Einstellungen geladen. Auch
% hier sind Sie frei diese direkt in der Praeambel zu laden.

% Folgende Pakete sind Minimalvoraussetzung.

% Anpassung von LaTeX an die deutsche Sprache. Zum Beispiel wird
% \chapter{} direkt als "`Kapitel"' uebersetzt usw. Ausserdem
% wird die Silbentrennung nach deutscher Rechtschreibung vorgenommen.
% Sonderzeichen, z.B. Umlaute, koennen direkt eingegeben werden.

\usepackage[ngerman]{babel}

% Kodierung fuer unixoide Systeme und Windowssysteme.

\usepackage[utf8]{inputenc}

% Schriftart Times New Roman.

\usepackage{times}

% Erweitert den Zeichenvorrat, so dass z.B. auch Umlaute im PDF-Dokument
% gefunden werden.

\usepackage[T1]{fontenc}

% Zur Einbindung von Bildern.

\usepackage{graphicx}

% Erweiterte enumerate-Umgebung.

\usepackage{enumerate}

% Verschiedene Pakete, die nuetzlich sind um Mathematik in LaTeX zu setzen.

\usepackage{amsmath, amssymb, amsthm, dsfont}

% Indentation

\usepackage{scrextend}

%%%%%%%%%%%%%%%%
% Seitenlayout %
%%%%%%%%%%%%%%%%

% DIV# gibt den Divisor für die Layoutberechnung an.
% Vergrößern des Divisors vergrößert den Textbereich.
% BCOR#cm gibt die Breite des Bundstegs an.
\usepackage[DIV14,BCOR2cm]{typearea}

% Abstand obere Blattkante zur Kopfzeile ist 2.54cm - 15mm
\setlength{\topmargin}{-15mm}

% Keine Einrueckung nach einem Absatz.

\parindent 0pt

% Abstand zwischen zwei Abs\"atzen.

\parskip 12pt

% Zeilenabstand.

\linespread{1.25}

% Inhaltsverzeichnis erstellen.

\usepackage{makeidx}
\makeindex


% Hier werden Makros und Umgebungen eingebunden. Diese werden separat
% in der Datei befehle.tex definiert. Sie sind frei diese Befehle auch
% direkt in der Praeambel zu definieren.
% Makros

\newcommand{\N}{\mathbb{N}} % natuerliche Zahlen
\newcommand{\Z}{\mathbb{Z}} % ganze Zahlen
\newcommand{\Q}{\mathbb{Q}} % rationale Zahlen
\newcommand{\R}{\mathbb{R}} % reelle Zahlen
\newcommand{\Rett}{Re}	% Schlechter Realteil		
\DeclareMathOperator{\Ret}{Re} % Guter Realteil

% Umgebungen für Definitionen, Sätze, usw.

\theoremstyle{definition}
%\newtheorem{defn}{Definition}[chapter]
\newtheorem{defn}{Definition}[section]
%\newtheorem{ex}{Beispiel}[chapter]

\theoremstyle{plain}
%\newtheorem{thm}{Satz}[chapter]
%\newtheorem{lem}{Lemma}[chapter]
%\newtheorem{cor}{Korollar}[chapter]
%\newtheorem{prop}{Proposition}[chapter]
\newtheorem{thm}{Satz}[section]
\newtheorem{lem}[thm]{Lemma}
\newtheorem{cor}[thm]{Korollar}
\newtheorem{prop}[thm]{Proposition}

\theoremstyle{remark}
\newtheorem*{bem}{Bemerkung}

\def\Satzrefname{Satz}


\begin{document}
% Auf der Titelseite und im Inhaltsverzeichnis sollen keine
% Seitenzahlen dargestellt werden.
\pagestyle{empty}

% Hier wird die Titelseite eingebunden.
\begin{titlepage}
  \begin{center}
    \vspace{1cm}
    \includegraphics[width=0.5\linewidth]{TU_Darmstadt_Logo.pdf}
    \vspace{12pt}
    
    \large{Fachbereich Mathematik}
    \vspace{2cm}
    
    \large{Proseminar zum Auswahlaxiom}
    \vspace{2cm}

    \huge{Topologie ohne Auswahlaxiom: Funktionenräume}
    \vspace*{2cm}    
    
		\large
    Fabian Gabel
    \vspace*{.5cm}

    22.01.2016 \\
    \vspace*{1cm}

    Betreuer: Dr. rer. nat. René Bartsch

    \vspace*{.5cm}

    %Zweiter Gutachter: Name des zweiten Gutachters\\[2ex]
    \tiny{Version vom \today}
  \end{center}
\end{titlepage}
\vspace*{\fill}


% Inhaltsverzeichnis erstellen.
\tableofcontents

% Ab sofort werden Seitenzahlen in der Kopfzeile angezeigt.
\pagestyle{headings}

% Einbinden der Einleitung.
%\chapter*{Einleitung}\index{Einleitung}
\addcontentsline{toc}{chapter}{Einleitung}\index{Einleitung}

\begin{addmargin}[2em]{2em}% 1em left, 2em right
  \textit{Zu jeder Menge $X$ von nicht leeren, zueinander disjunkten Mengen gibt es eine Menge, die von jedem Elemnt von $X$ genau ein Element enthält.} 
  \flushright(Auswahlaxiom)
\end{addmargin}

Das Auswahlaxiom \AC ist eines der kontroversesten Axiome der Mengenlehre.
Es wurde schon 1904 nach der ersten expliziten Formulierung durch \textsc{Zermelo} nur zögernd aufgenommen. 
Ein Grund dafür liegt wohl darin, dass sich aus der anscheinlich intuitiven Vorstellung, die man mit \AC verbinden kann, infolge des inkonstruktiven Charakters dieses Axioms oftmals unintuitive Aussagen folgern lassen.
Bekannte Beispiele hierfür liefern der Wohlordnungssatz oder das \textsc{Banach}-\textsc{Tarski}-Paradoxon.
Trotzdem bildet \AC einen wichtigen Bestandteil mathematischer Theorien, wie auch der Theorie topologischer Räume.
In dieser Arbeit stehen basierend auf \cite{herrlich2006axiom} speziell topologische Funktionenräume im Vordergrund unter dem Aspekt der Gültigkeit verschiedener Versionen von \textsc{Áscoli}-Sätzen.

Es werden Grundlagen, wie sie in Rahmen einer Vorlesung zu allgemeiner Topologie vermittelt werden, vorausgesetzt.
Sämtliche unbewiesene Aussagen dieses Typus lassen sich jedoch in der vorhandenen Lehrbuchliteratur auffinden.
Für diesen Text wurden hierzu insbesondere \cite{bartsch2015allgemeine},\cite{ebbinghaus2003einfuhrung},\cite{preuss1972topologische} und \cite{kelley1975general} herangezogen.

Im Einführungskapitel dieser Arbeit werden grundlegende Aussagen der Topologie und Mengenlehre bewiesen, die es ermöglichen den Charakter ihrer Beziehung zu \AC und dem \textsc{Boole}schen Primidealsatz (\PIT) zu beschreiben. 
Es werden unterschiedliche, im Rahnmen von \ZF zunächst voneinander unabhängige Kompaktheitsbegriffe auf topologischen Räumen eingeführt und anschließend in Relation zueinander gesetzt.
Als zentrale Aussage dieses Kapitels steht ein Satz, welcher eine Abschwächung von \AC, namlich den \textsc{Boole}schen Primidealsatz \PIT, mit Kompaktheitsaussagen über Produkttopologien verbindet.

Das zweite Kapitel behandelt die Topologie von Funktionenräumen im Hinblick auf \textsc{Áscoli}-Sätze.
Es werden unterschiedliche Formulierungen dieser Kompaktheitsaussagen vorgestellt und deren Äquivalenz zu Axiomen der Mengenlehre bewiesen.
Insbesondere wird aufgezeigt inwieweit sich \textsc{Áscoli}-Sätze im Rahmen von \ZF beweisen lassen, wenn man alternative Kompaktheitsbegriffe zugrunde legt.

Abschließend fasst ein Fazit die wichtigsten Ergebnisse zusammen.



% Einbinden der einzelnen Kapitel (zur Übersicht eine Datei je Kapitel)

\chapter{Mengentheoretische und topologische Grundlagen}

%\subsection{Das Auswahlaxiom}

%\subsection{Der \textsc{Boole}sche Primidealsatz}

\section{Kompaktheitsbegriffe auf topologischen Räumen}

Das Auswahlaxiom oder Abschwächungen davon dienen oftmals als Bindeglied zwischen unterschiedlichen Konzepten in der Mathematik, indem sie ihre Äquivalenz zeigen. 
Ein für die vorliegende Ausarbeitung zentrales Konzept ist das der Kompaktheit topologischer Räume. 
Es exisiteren unterschiedliche Kompaktheitsbegriffe, die im Rahmen von \textbf{ZF} nicht notwendig äquivalent sind.

\begin{defn}
  Sei $(X,\tau)$ ein topologischer Raum.
  Ein Punkt $x \in X$ heißt Adhärenzpunkt eines Filters $\varphi \in \mathcal{F}(X)$, falls es einen gegen $x$ bezüglich $\tau$ konvergenten Oberfilter $\psi \supseteq \varphi$ gibt.
\end{defn}

\begin{prop}
  \label{prop:adherence}
  Sei $(X,\tau)$ ein topologischer Raum.
  Ein Punkt $x \in X$ ist genau dann Adhärenzpunkt eines Filters $\varphi \in \mathcal{F}(X)$, falls
  \begin{displaymath}
    x \in \bigcap_{F \in \varphi} \overline{F}.
  \end{displaymath}
\end{prop}

\begin{proof}
  Sei zunächst $x$ Adhärenzpunkt eines Filters $\varphi$ auf $X$.
  Nach Definition existiert ein gegen $x$ konvergenter Oberfilter $\psi \supseteq \varphi$, also gilt $\psi \supseteq \dot x \cap \tau$
  Somit gilt aufgrund der Filtereigenschaft von $\psi$ für alle $F \in \varphi$ und alle $U \in x \cap \tau$, dass  $F \cap U \neq \emptyset$.
  Also besitzt jede Umgebung $U$ von $x$ mit allen $F \in \varphi$ einen nichtleeren Schnitt, was bedeutet, dass $x \in \overline F$ für alle $F  \in \varphi$

  Sei andererseits $x \in \bigcap_{F \in \varphi} \overline F$.
  Nach Definition des Abschlusses einer Menge, gilt somit für alle $U \in \dot x \cap \tau$, dass $F \cap U \neq \emptyset$.
  Es ist also $\varphi \cup (\dot x \cap \tau)$ eine Filtersubbasis.
  Man bezeichne mit $\psi$ den davon erzeugten Filter.
  Nach Konstruktion gelten $\psi \supseteq \varphi$ und $\psi \overset{\tau}{\to} x$.

\end{proof}

\begin{defn}
  Ein topologischer Raum $(X,\tau)$ heißt
  \begin{enumerate}[(1)]
    \item kompakt, \textit{falls sich aus jeder offenen Überdeckung von $X$ eine Teilüberdeckung auswählen lässt},
    \item filterkompakt, \textit{falls jeder Filter auf $X$ einen Häufungspunkt besitzt},
    \item ultrafilterkompakt, \textit{falls jeder Ultrafilter auf $X$ konvergiert},
    \item \textsc{Alexandroff}-\textsc{Urysohn}-kompakt, \textit{falls jede unendliche Teilmenge von $X$ einen vollständigen Häufungspunkt besitzt},
    \item \textsc{Tychonoff}-kompakt, \textit{falls $X$ homöomorph zu einem abgeschlossenen Teilraum eines Hilbert-Würfels $[0,1]^I$ ist}.
  \end{enumerate}
\end{defn}


\begin{thm}
\label{thm:compactness}
  Sei $(X,\tau)$ ein topologischer Raum. Dann gilt:
  \begin{enumerate}[(1)]
    \item Der Raum $X$ ist genau dann kompakt, wenn er filterkompakt ist. 
    \item Ist $X$ filterkompakt, dann auch ultrafilterkompakt.
  \end{enumerate}
\end{thm}

\begin{proof}
  $(1)$: 
  Sei $(X,\tau)$ kompakt und $\varphi$ ein Filter auf $X$.
  Angenommen, $\varphi$ besitze keinen Häufungspunkt.
  Dann gilt nach Proposition \ref{prop:adherence}
  \begin{displaymath}
    \bigcap_{F \in \varphi} \overline{F} = \emptyset.
  \end{displaymath}
  Folglich ist 
  \begin{displaymath}
    X = \bigcup_{F \in \varphi} X \setminus \overline{F} \subseteq \bigcup_{F \in \varphi} X \setminus F 
  \end{displaymath}
  eine offene Überdeckung von $X$.
  Da $X$ als kompakt vorausgesetzt wird, existieren $F_1,\dots,F_n \in \varphi$ mit
  \begin{displaymath}
    \bigcup_{i = 1}^n X \setminus F_i = X
  \end{displaymath}
  woraus wiederum
  \begin{displaymath}
    X \setminus \bigcup_{i = 1}^n F_i = \bigcap_{i = 1}^n F_i = \emptyset
  \end{displaymath}
  folgt. 
  Dies steht jedoch im Widerspruch zur Filtereigenschaft.

  Sei umgekehrt $X$ filterkompakt und $(O_i)_{i \in I}$ eine offene Überdeckung von $X$.
  Angenommen es existiere keine Teilüberdeckung.
  Man betrachte nun die Familie
  \begin{displaymath}
    \mathfrak{B} := \left\{ X \setminus \mathcal{O} \mid \mathcal{O} = \bigcup_{k = 1}^n O_{i_k}\right\}.
  \end{displaymath}
  Diese Familie ist nichtler, da $X \setminus O_i \in \mathcal{O}$ für alle $i \in I$ gilt und zudem abgeschlossen unter endlichen Schnitten.

  Folglich ist $\mathfrak{B}$ eine Filterbasis und es bezeichne $\varphi$ den durch Obermengenbildung davon erzeugten Filter.
  Nach Voraussetzung existiert ein Oberfilter $\psi \supseteq \varphi$ der bezüglich der Topologie $\tau$ gegen ein $x \in X$ konvergiert.
  Aus der Überdeckungseigenschaft folgt, dass zusätzlich ein $i_0 \in I$ existiert mit $x \in O_i$.
  Aufgrund der Konvergenz des Filters $\psi$ folgt $O_i \in \psi$.
  Dies steht jedoch im Widerspruch zu $\psi \supseteq \mathfrak{B}$, denn deswegen gilt bereits $X \setminus O_i \in \psi$.

  $(2)$: Da Ultrafilter bereits bezüglich Inklusion maximal sind, also keine echten Oberfilter besitzen, impliziert filterkompaktheit bereits ultrafilterkompaktheit.
\end{proof}

\begin{defn}
    \UFT, der Ultrafiltersatz: 
      \begin{addmargin}[2em]{2em}% 1em left, 2em right
        \textit{Jeder Filter auf einer Menge lässt sich zu einem Ultrafilter erweitern.}
      \end{addmargin}
\end{defn}

\begin{thm}
  Es sind äquivalent:
  \begin{enumerate}
    \item Ein topologischer Raum ist genau dann kompakt, wenn er ultrafilterkompakt ist.
    \item \UFT.
  \end{enumerate}
\end{thm}

\begin{thm}
  Es sind äquivalent:
  \begin{enumerate}[(1)]
    \item Ein topologischer Raum ist genau dann kompakt, wenn er \textsc{Alexandroff}-\textsc{Urysohn}-kompakt ist.
    \item Ein topologischer Raum ist genau dann ultrafilterkompakt, wenn er \textsc{Alexandroff}-\textsc{Urysohn}-kompakt ist.
    \item \UFT.
  \end{enumerate}
\end{thm}

\section{Der \textsc{Boole}sche Primidealsatz}

\begin{defn}
 \item \AC, das \textit{Auswahlaxiom}.
 \item Lemma von \textsc{Zorn}.
\end{defn}

\begin{defn}
    \PIT, der \textit{\textsc{Boole}sche Primidealsatz}:
      \begin{addmargin}[2em]{2em}% 1em left, 2em right
        \textit{Jede \textsc{Boole}sche Algebra besitzt ein maximales Ideal.}
      \end{addmargin}
\end{defn}

\begin{thm}
  \label{thm:acpit}
  \AC impliziert \PIT.
\end{thm}

\begin{proof}
  Sei $X$ eine \textsc{Boole}sche Algebra und $I$ ein Ideal. 
  Man betrachte die Menge
  \begin{displaymath}
    \mathcal{A} := \{J \subsetneq X \mid \text{ $J$ ist ein $I$ umfassendes Ideal } \}.
  \end{displaymath}
  Um das \textsc{Zorn}sche Lemma anwenden zu können, betrachtet man nun eine Kette $\mathcal{K} \subseteq \mathcal{A}$.
  Es sei 
  \begin{displaymath}
    B := \bigcup_{J \in \mathcal{K}} J.
  \end{displaymath}

  Es soll nun gezeigt werden, dass $B \in \mathcal{A}$ gilt.
  Da nach Voraussetzung $\mathcal{K}$ nur aus $I$ umfassenden Idealen besteht, gilt auch $I \subseteq B$.
  Sind $x,y \in B$ so existieren $J_x, J_y \in \mathcal{K}$ mit $x \in J_x$ und $y \in J_y$.
  Da $\mathcal{K}$ eine Kette ist, kann man aus Symmetriegründen ohne Beschränkung der Allgemeinheit annehmen $J_x \subseteq J_y$ und damit $x \in J_y$.
  Da $J_y$ nach Voraussetzung ein Ideal ist gilt $x + y \in J_y \subseteq B$.
  Ist andererseits $x \in B$ und $a \in X$ mit $a \leq x$, so existiert wie zuvor ein $J \in \mathcal{K}$ mit $x \in J$ und, da $J$ ein Ideal ist, gilt auch $a \in J$ und damit $a \in B$.
  Also ist $B$ ein Ideal, welches $I$ umfasst, das bedeutet $B \in \mathcal{A}$.

  Nach Konstruktion gilt $J \subseteq B$ für alle $J \in \mathcal{K}$, und damit ist $B$ eine obere Schranke von $\mathcal{K}$ in $\mathcal{A}$.
  Das \textsc{Zorn}sche Lemma garantiert nun die Existenz maximaler Elemente in $\mathcal{A}$, also die Existenz eines maximalen, $I$ umfassenden Ideals.
\end{proof}

\begin{bem}
  Auf den ersten Blick scheint der bewiesene Satz \ref{thm:acpit} stärker zu sein als \PIT.
  Es hätte bereits genügt, anstatt $\mathcal{A}$ die Menge aller echten Ideale in $X$ zu betrachten.
  Tatsächlich sind die beiden Aussagen jedoch äquivalent.
\end{bem}

\begin{proof}
  
  Um zu zeigen, dass aus \PIT bereits zu jedem Ideal die Existenz maximaler umfassender Ideale folgt, betrachte man eine \textsc{Boole}sche Algebra $X$ mit Ideal $I$.
  
  Es induziert $I$ folgendermaßen eine Äquivalenzrelation auf $X$.
  Es soll für $x,y \in X$ gelten, dass 
  \begin{displaymath}
    x \sim y \quad \text{genau dann, wenn} \quad (x \cdot -y) + (y \cdot -x) \in I
  \end{displaymath}
  gilt.
  Auf der Menge der Restklassen 
  \begin{displaymath}
    X / I := \{[x] \mid x \in X\}
  \end{displaymath}
  lassen sich nun über die Repräsentanten folgende Operationen definieren:
  \begin{displaymath}
    [x] + [y] := [x+y], \quad
    [x] \cdot [y] := [x \cdot y], \quad
    -[x] := [-x]
  \end{displaymath}
  Damit wird $X/I$ zu einer \textsc{Boole}schen Algebra.
  Setzt man \PIT voraus, so existiert ein maximales Ideal $K \subseteq X/I$
  Betrachtet man nun die Menge
  \begin{displaymath}
    J := \{x \in X \mid [x] \in K\},
  \end{displaymath}
  so zeigt man nun, dass $J$ ein maximales Ideal in $X$ ist, welches $I$ umfasst.
  
  Dass $J$ ein Ideal ist, folgt sofort aus der Definition der Verknüpfungsoperationen auf $X/I$ und der Tatsache, dass $K$ ein Ideal ist.
  Für alle $x \in I$ gilt zudem $[x] = [0] \in X/I$. 
  Da $K$ als Ideal definitionsgemäß $[0]$ enthält, folgt $[x] \in K$, also $I \subseteq J$.
  
  Angenommen $J$ sei nicht maximal. Sei $J' \subsetneq X$ ein $J$ umfassendes Ideal.
  Dann gilt 
  \begin{displaymath}
    K \subseteq J'/I := \{[x] \mid x \in J\},
  \end{displaymath}
  denn für $[x] \in K$ ist $x \in J \subseteq J'$ und damit $[x] \in J'/X$.
  Aufgrund der Maximalität von $K$ folgt sofort $K := J/I$ und damit $J' = J$.

  Also ist $J$ maximales $I$ umfassendes Ideal.
\end{proof}

In \textsc{Boole}schen Algebren existiert das zum Ideal duale Konzept des Filters. 

\begin{defn}
  Filter
\end{defn}

Es lässt sich somit eine zu \PIT duale und damit äquivalente Formulierung für Filter finden:
  \begin{addmargin}[2em]{2em}% 1em left, 2em right
    \textit{Jede \textsc{Boole}sche Algebra besitzt einen maximalen Filter.}
  \end{addmargin}

Dass \PIT schwächer ist als \AC, kann man in QUELLE nachlesen.
Die Bedeutung von \PIT ist jedoch nicht zu unterschätzen QUELLE

\begin{thm}
  Äquivalent sind
\begin{enumerate}[(1)]
    \item \PIT.
    \item \UFT.
    \item \textsc{Tychonoff}-Satz für \textsc{Hausdorff}-Räume: 
      \begin{addmargin}[2em]{2em}% 1em left, 2em right
        Produkte kompakter \textsc{Hausdorff}-Räume sind kompakt.
      \end{addmargin}
    \item \textsc{Tychonoff}-Satz für endliche diskrete Räume: 
      \begin{addmargin}[2em]{2em}% 1em left, 2em right
        Produkte endlicher diskreter Räume sind kompakt.
      \end{addmargin}
    \item \textsc{Hilbert}-Würfel $[0,1]^I$ sind kompakt.
    \item \textsc{Kantor}-Würfel $\mathbf{2}^I$ sind kompakt.
  \end{enumerate}

  \begin{proof}
    $(1)\Rightarrow(2)$: 
    Interpretiert man \PIT im Sinne von Filtern, so folgt dies sofort, da jede Potenzmengenalgebra eine \textsc{Boole}sche Algebra ist.

    $(2)\Rightarrow(3)$:

  \end{proof}

\end{thm}


%\section{Der \textsc{Ascoli}-Satz und das Auswahlaxiom}
\chapter{Der \textsc{Ascoli}-Satz und das Auswahlaxiom}

\textsc{Ascoli}-Sätze gelten bezeichnend für Aussagen, die es gestatten Kompaktheit von Teilmengen von Funktionenräumen bezüglich starker Topologien aus der Kompaktheit bezüglich einer schwächeren Topologie unter zusätzlichen Annahmen zu folgern. 
Dabei existieren unterschiedliche Formulierungen.

Dieser Abschnitt ist eine ausgestaltete Darstellung des Kapitels  
\begin{addmargin}[2em]{2em}% 1em left, 2em right
  \textit{Disasters in Topology III: Function Spaces (The Ascoli Theorem)}
\end{addmargin}
aus \cite{herrlich2006axiom}.
Weitere allgemeinere Versionen des \textsc{Ascoli}\hyp{}Satzes finden sich in \cite{bartsch2015allgemeine} oder \cite{poppe1974compactness}.

Ziel dieses Kapitels ist es zu untersuchen inwieweit sich \textsc{Ascoli}\hyp{}Sätze im Rahmen von \ZF beweisen lassen und welche Erweiterungen des Axiomensystems \ZF gegebenenfalls für den Beweis benötigt werden.
Eine wichtige Rolle für diese Beweise wird der \textsc{Boole}sche Primidealsatz \PIT spielen, dessen Gültigkeit, wie im vorangehenden Kapitel gezeigt, zu unterschiedlichen Kompaktheitsaussagen äquivalent ist.

\section{Der topologische \textsc{Ascoli}-Satz}

In diesem Abschnitt soll eine allgemeine topologische Version des \textsc{Ascoli}\hyp{}Satzes behandelt werden.

\begin{defn}
  Es sei $(X,\tau)$ ein topologischer Raum und $(Y,d)$ ein metrischer Raum versehen mit der durch die Metrik induzierte Topologie $\tau_d$. Eine Teilmenge $F \subseteq Y^X$ heißt \textit{gleichgradig stetig}, wenn gilt:
  \begin{addmargin}[2em]{2em}% 1em left, 2em right
    Für alle $x \in X$ und $\varepsilon > 0$ existiert ein $V \in \dot{x} \cap \tau$, sodass \\
    für alle $f \in F$ gilt, dass $f(V) \subseteq U_\varepsilon(f(x))$.
  \end{addmargin}
\end{defn}

Für zwei topologische Räume $(X,\tau)$ und $(Y,\sigma)$ bezeichne im Folgenden $C_{\co}(X,Y)$ den topologischen Funktionenraum $(C(X,Y),\tau_{\co})$ der stetigen Funktionen von $X$ nach $Y$ ausgestattet mit der \textit{kompakt\hyp{}offenen} Topologie.
Diese Topologie wird erzeugt von Subbasiselementen der Form
\begin{displaymath}
  (K,O) := \{f \in C(X,Y) \mid K \subseteq X \text{ ist kompakt und } O \in \sigma\}.
\end{displaymath}

\begin{defn}[Topologischer \textsc{Ascoli}-Satz]
  Sei $(X,\tau)$ ein lokalkompakter \textsc{Hausdorff}\hyp{}Raum und $(Y,d)$ ein metrischer Raum.
  Für jeden Teilraum $(F,\tau_{\co})$ von $C_{\co}(X,Y)$ sind äquivalent:
  \begin{enumerate}
    \item[(a)] $F$ ist $\tau_{\co}$\hyp{}kompakt.
    \item[(b)] \begin{enumerate}
        \item[($\alpha$)] Für alle $x \in X$ ist die Menge $F(x) = \{f(x) \mid f \in F \}$ kompakt in $Y$.
        \item[($\beta$)]  $F$ ist abgeschlossen in $Y^X$ bezüglich der Topologie der punktweisen Konvergenz $\tau_p$.
        \item[($\gamma$)] $F$ ist gleichgradig stetig auf $X$.
      \end{enumerate}
  \end{enumerate}
\end{defn}

Die Hauptaussage des Satzes ist also die Folgerung der Kompaktheit einer Menge stetiger Funktionen bezüglich einer starken Topologie, in diesem Fall der kompakt\hyp{}offenen Topologie, aus der Kompaktheit bezüglich der schwächeren punktweisen Topologie.
Eine verbreitete Variation des Satzes verzichtet auf Teil (b,$\beta$) und ersetzt die Kompaktheit in (a) durch relative Kompaktheit.
Man beachte, dass die auf dem Teilraum $F$ verwendete Topologie $\tau_{\co}$ die Spur der kompakt\hyp{}offenen Topologie auf $C(X,Y)$ ist.
Es soll nun gezeigt werden, dass die Gültigkeit dieses Satzes äquivalent zum \textsc{Boole}schen Primidealsatz (\PIT) ist.

\begin{thm}
  \label{thm:ascoliPIT}
  Äquivalent sind
  \begin{enumerate}
    \item[(1)] Der topologische \textsc{Ascoli}-Satz.
    \item[(2)] \PIT.
  \end{enumerate}
\end{thm}

\begin{proof}
  (1) $\Rightarrow$ (2): Sei $X$ eine Menge und $\mathbf{2}$ die zweielementige Menge $\{0,1\}$. Beide Mengen seien  mit der diskreten Topologie ausgestattet. 
  Die diskrete Topologie lässt sich auch als die von der \textit{diskreten Metrik} induzierte Topologie auffassen, somit kann $\mathbf{2}$ insbesondere als metrischer Raum betrachtet werden. 

  Es sei $F := C(X, \mathbf{2})$.
  Dann gilt $C(X, \mathbf{2}) = \mathbf{2}^X$, da $X$ nach Voraussetzung diskret ist.
  Es folgt, dass $\mathbf{2} = F(x)$ für alle $x \in X$ als endlicher Raum kompakt ist.
  Nach Definition ist $\mathbf{2}^X$ als Gesamtraum abgeschlossen bezüglich jeder Topologie, also insbesondere auch bezüglich der punktweisen Topologie $\tau_p$.
  Zuletzt ist $F$ auch gleichgradig stetig, da $X$ diskret und somit Einpunktmengen offen sind und man daher für alle $x \in X$ und $\varepsilon > 0$ die Menge $\{x\}$ als offene Umgebung von $x$ wählen kann.
  Dann gilt nämlich für alle $f \in F$, aufgrund der Definitheit der Metrik für  $d(f(x),f(x)) = 0 < \varepsilon$.
  Damit sind alle Bedingungen aus $(b)$ des \textsc{Ascoli}\hyp{}Satzes erfüllt und es folgt, dass $\mathbf{2}^X$ kompakt bezüglich der kompakt\hyp{}offenen Topologie ist.
  Da die kompakt\hyp{}offene Topologie stärker als die punktweise Topologie ist, folgt damit die Kompaktheit von $\mathbf{2}^X$ bezüglich der punktweisen Topologie.
  Dies ist nach Satz \ref{thm:pitequivalence} äquivalent zu \PIT.

  (2) $\Rightarrow$ (1): Es seien $(X,\tau)$ ein lokalkompakter \textsc{Hausdorff}\hyp{}Raum und $(Y,d)$ ein metrischer Raum.
  Des Weiteren sei $F$ ein Teilraum vom $C_{\co}(X,Y)$.

  (a) $\Rightarrow$ (b, $\alpha$): 
  Nach Voraussetzung ist $F$ ein bezüglich der kompakt\hyp{}offenen Topologie $\tau_{\co}$ kompakter Teilraum von $C(X,Y)$.
  Somit ist er insbesondere bezüglich der Spur der schwächeren Topologie der punktweisen Konvergenz $\tau_p$ ein kompakter Teilraum des \textsc{Hausdorff}\hyp{}Raumes $Y^X$.
  Die Topologie der punktweisen Konvergenz lässt sich als Initialtopologie bezüglich der für alle $x \in X$ definierten kanonischen Projektionen 
  \begin{displaymath}
    \pi_x : Y^X \to Y, \quad f \mapsto f(x)
  \end{displaymath}
  beschreiben.
  Als stetiges Bild eines Kompaktums ist somit auch $\pi_x(F) = F(x)$ kompakt in $(Y,\tau_d)$.

  (a) $\Rightarrow$ (b, $\beta$):  
  Als kompakter Teilraum eines \textsc{Hausdorff}\hyp{}Raumes ist $F$ bezüglich $\tau_p$ auch abgeschlossen in $Y^X$.

  (a) $\Rightarrow$ (b, $\gamma$):
  Es soll gezeigt werden, dass $F$ gleichgradig stetig ist.
  Sei dazu $x \in X$ und $\varepsilon > 0$.
  Für jedes $f \in F$ ist die Menge
  \begin{displaymath}
    B_f := B_{\frac{\varepsilon}{2}} = \left\{ y \in Y \mid d(f(x),y) < \tfrac{\varepsilon}{2}\right\} 
  \end{displaymath}
  offen bezüglich der von der Metrik $d$ auf $Y$ induzierten Topologie.
  Da $F \subseteq C(X,Y)$, ist das Urbild $f^{-1}(B_f)$ ebenfalls offen bezüglich $\tau$.
  Nach Voraussetzung ist $X$ lokalkompakt.
  Daher existiert eine kompakte Umgebung $K_f$ von $x$ mit $K_f \subseteq f^{-1}(B_f)$.
  Damit gilt $f(K_f) \subseteq B_f$ und es definiert
  \begin{displaymath}
    U_f := F \cap (K_f, B_f) = \left\{g \in F \mid g(K_f) \subseteq B_f \right\} \tag{$\ast$}
  \end{displaymath}
  eine bezüglich der Teilraumtopologie auf $F$ offene Umgebung von $f$.
  Man betrachte nun die Evaluationsabbildung 
  \begin{displaymath}
    \omega: X \times F \to Y, \quad (y,g) \mapsto g(y).
  \end{displaymath}
  Unter Berücksichtigung von ($\ast$) folgt $\omega(K_f \times U_f) \subseteq B_f$.
  Es definiere $\mathcal{C}$ die Menge aller Tripel $(f,K,U)$, wobei $f \in F$, $K$ eine kompakte Umgebung von $x$ in $X$ und $U$ eine offene Umgebung von $f$ in $F$ mit $\omega(K \times U) \subseteq B_f$ sei.
  Dann liefert
  \begin{displaymath}
    \mathfrak{U} := \{ U \subseteq F \mid \text{es existieren $f \in F$, $K \subseteq X$ mit $(f,K,U) \in \mathcal C$} \}
  \end{displaymath}
  eine offene Überdeckung von $F$, da für alle $f \in F$ nach Konstruktion $(f,K_f,U_f) \in \mathcal{C}$, also auch $U_f \in \mathfrak{U}$ gilt.
  Nach Voraussetzung ist $F$ kompakt, daher existiert eine endliche Teilüberdeckung durch $U_1,\dots,U_n \in \mathfrak{U}$.
  Für alle $i \in \{ 1,\dots,n\}$ wähle man nun $f_i \in F$ und $K_i \subseteq X$ mit $(f_i, K_i, U_i) \in \mathcal{C}$.
  Des Weiteren definiere man 
  \begin{displaymath}
    U := \bigcap_{i = 1}^n K_i.
  \end{displaymath}
  Es ist $U$ eine nicht notwendig offene aber als endlicher Durchschnitt kompakter Mengen eine kompakte Umgebung von $x \in X$.
  Nun soll gezeigt werden, dass für alle $f \in F$ folgt, dass $f(U) \subseteq U_\varepsilon(f(x))$ gilt.
  Da $X$ nach Voraussetzung lokalkompakt ist, enthält $U$ folglich auch eine offene Umgebung $V$ von $x$ für die dann entsprechend $f(V) \subseteq U_\varepsilon(f(x))$ gilt.
  Damit ist dann $F$ gleichgradig stetig.

  Für jedes $f \in F$ existiert aufgrund der Überdeckungseigenschaft ein $i \in \{ 1,\dots,n\}$ mit $f \in U_i$.
  Sei des Weiteren $y \in U$ gegeben.
  Dann gilt
  \begin{displaymath}
    f(y) = \omega(y,f) \in \omega(U \times U_i) \subseteq \omega(K_i \times U_i) \subseteq B_{f_i},
  \end{displaymath}
  also $d(f_i(x),f(y)) < \tfrac{\varepsilon}{2}$.
  Insbesondere gilt für $x \in U$, da nach Voraussetzung $x \in K_i$ für alle $i \in \{ 1,\dots,n\}$, dass $d(f_i(x),f(x)) < \tfrac{\varepsilon}{2}$.
  Unter Verwendung der Dreiecksungleichung folgt letztlich
  \begin{displaymath}
    d(f(x),f(y)) \leq d(f(x),f_i(x)) + d(f_i(x),f(y)) < \frac{\varepsilon}{2} + \frac{\varepsilon}{2} = \varepsilon
  \end{displaymath}
  und damit auch $f(U) \subseteq U_\varepsilon(f(x))$.
  Folglich ist $F$ gleichgradig stetig.

  (b) $\Rightarrow$ (a)
  Nach Voraussetzung ist $F(x)$ als Teilraum von $Y$ kompakt und besitzt zudem als metrischer Raum die \textsc{Hausdorff}\hyp{}Eigenschaft.
  %%%%%%%% PIT %%%%%%%%
  Der \textsc{Tychonoff}\hyp{}Satz für \textsc{Hausdorff}\hyp{}Räume \ref{thm:pitequivalence}(3) besagt, dass dann auch $\prod_{x \in X} F(x) \subseteq Y^X$ kompakt ist.
  %%%%%%%%%%%%%%%%%%%%%
  Die Voraussetzung, dass $F$ abgeschlossen in $Y^X$ bezüglich der Produkttopologie $\tau_p$ ist, liefert, dass $F$ auch abgeschlossen im Teilraum $\prod_{x \in X} F(x)$ ist, da $F \subseteq \prod_{x \in X} F(x)$ gilt.
  Als abgeschlossene Teilmenge eines Kompaktums ist somit $F$ auch kompakt bezüglich der Topologie der punktweisen Konvergenz.

  Es fehlt zu zeigen, dass $F$ ebenfalls bezüglich der kompakt\hyp{}offenen Topologie $\tau_{\co}$ kompakt ist.
  Hierfür soll die Inklusion $\tau_{\co} \subseteq \tau_p$ nachgewiesen werden.
  Dazu betrachte man ein Subbasiselement 
  \begin{displaymath}
    V := (K,U) \cap F \in \tau_{\co}
  \end{displaymath}
  der kompakt\hyp{}offenen Topologie auf dem Teilraum $F$ und ein $f \in V$.
  Es gilt nun zu zeigen, dass $V$ auch in $\tau_p$ offen ist.
  Da $f(K) \subseteq U$ und $U \in \sigma$ gilt, folgt für alle $x \in K$
  \begin{displaymath}
    r_x := \inf\{d(f(x),y) \mid y \in (Y \setminus U) \} > 0.
  \end{displaymath}
  Dann ist 
  \begin{displaymath}
    U_x := \{ z \in X \mid d(f(x),f(z)) < \tfrac{r_x}{2}\} = f^{-1}\left( U_{\tfrac{r_x}{2}}(f(x))\right)
  \end{displaymath} 
  eine offene Umgebung von $x \in X$ aufgrund der Stetigkeit von $f$.
% Denn für alle $z \in U_x$ mit $\varepsilon := \tfrac{r_x}{2} - d(f(x),f(z)) > 0$ gilt, dass aufgrund der Stetigkeit von $f$ eine Umgebung $U' \in \tau$ von $z$ existiert, sodass $f(U') \subseteq U_\varepsilon(f(z))$ gilt.
% Daraus folgt für alle $u \in U'$ mit der Dreiecksungleichung
% \begin{align*}
%   d(f(x),f(u)) &\leq d(f(x),f(z)) + d(f(z),f(u)) \\
%                &< \varepsilon + d(f(z),f(u)) = \frac{r_x}{2} - d(f(x),f(z)) + d(f(x),f(z)) =  \frac{r_x}{2}
% \end{align*}
% und damit schließlich $U' \subseteq U_x$, also die Offenheit von $U_x$ in $X$.
  Zudem bildet $\mathfrak{U} := \{U_x \mid x \in K\}$ eine offene Überdeckung von $K$. 
  Nach Voraussetzung ist $K$ kompakt, daher existieren $x_1,\dots,x_n \in K$, sodass $K \subseteq \bigcup_{i=1}^n U_{x_i}$.
  Es sei $r:= \min\{r_{x_1},\dots,r_{x_n}\}$. 
  Damit ergibt sich für alle $x \in K$ und alle $y \in (Y \setminus U)$ die Ungleichung $d(f(x),y) \geq \tfrac{r}{2}$.
  Ist andererseits $x \in K$ und $d(f(x),y) < \tfrac{r}{2}$, so folgt daraus $y \in U$.

  Da $F$ nach Voraussetzung gleichgradig stetig auf $X$ ist, existiert insbesondere für alle $x \in K$ eine offene Umgebung $W$ von $x$, sodass 
  \begin{displaymath}
    d(g(x),g(z)) < \tfrac{r}{4}, \quad \text{für alle $g \in F$ und $z \in W$}.
  \end{displaymath}
  Betrachtet man nun die auf obige Weise erzeugte Menge $\mathcal{C}$ aller geordneten Paare $(x,W)$, so bildet
  \begin{displaymath}
    \mathfrak{W} := \{W \subseteq X \mid \text{Es existiert ein $x \in K$ mit $(x,W) \in \mathcal{C}$} \}
  \end{displaymath}
  eine offene Überdeckung von $K$.
  Aufgrund der Kompaktheit von $K$ lässt sich eine endliche Teilüberdeckung $W_1,\dots,W_m \in \mathfrak{W}$ auswählen.
  Für alle $i \in \{ 1,\dots,m\}$ wähle man entsprechend der Definition $\mathfrak{W}$ ein $x_i$ mit $(x_i,W_i) \in \mathcal{C}$.
  Dann ist
  \begin{displaymath}
    B_f := \{ g \in F \mid d(f(x_i),g(x_i)) < \tfrac{r}{4} \text{ für } i \in \{ 1,\dots,m\}\}
  \end{displaymath}
  als Basisumgebung von $\tau_p$ insbesondere eine offene Umgebung von $f$.

  Es soll nun gezeigt werden, dass $B_f \subseteq V$ gilt.
  Dann ist nämlich $V = \bigcup_{f \in V} B_f$ als Vereinigung offener Mengen offen bezüglich $\tau_p$ und die Behauptung $\tau_{\co} \subseteq \tau_p$ folgt, da $V$ beliebig war.
  Sei dazu also $g \in B_f$.
  Für jedes $x \in K$ existiert ein $i \in \{1,\dots,m\}$ mit $x \in W_i$.
  Nach Konstruktion der $W_i$ gilt $d(g(x_i),g(x)) < \tfrac{r}{4}$.
  Da $g \in B_f$ gilt zudem die Ungleichung $d(f(x_i),g(x_i)) < \tfrac{r}{4}$.
  Daraus folgert man
  \begin{displaymath}
    d(f(x_i),g(x)) \leq d(f(x_i),g(x_i)) + d(g(x_i),g(x)) < \frac{r}{2},
  \end{displaymath}
  was wiederum $g(x) \in U$ und damit $g(K) \subseteq U$ impliziert. 
  Also ist $g \in (K,U)$, woraus $g \in V$ folgt.
  Damit ist der Beweis vollständig.
\end{proof}

\section{Alternative Kompaktheitsbegriffe im \textsc{Ascoli}-Satz}

Betrachtet man den Beweis von Satz \ref{thm:ascoliPIT}, so fällt auf, dass die Implikation ((a) $\Rightarrow$ (b)) in Satz \ref{thm:ascoliPIT} bereits im Rahmen von \ZF gültig ist.
Der Beweis dieses Satzes verwendet lediglich einer Stelle der Beweisrichtung ((b) $\Rightarrow$ (a)) den laut Satz \ref{thm:pitequivalence} zu \PIT äquivalenten \textsc{Tychonoff}\hyp{}Satz für kompakte \textsc{Hausdorff}\hyp{}Räume.
Unter der Verwendung des Kompaktheitsbegriffes der Ultrafilterkompaktheit ist diese Schlussweise gemäß Lemma \ref{lem:uftprod} jedoch bereits in \ZF möglich.  
Dies motiviert den Versuch zu untersuchen, inwiefern der topologische \textsc{Ascoli}\hyp{}Satz unter Verwendung anderer Kompaktheitsbegriffe seine Gültigkeit behält und ob sich diese Versionen des \textsc{Ascoli}\hyp{}Satzes bereits in \ZF beweisen lassen.

\begin{lem}
  \label{lem:filterclusterpoint}
  Es sei $(X,\tau)$ ein topologischer Raum, sodass die darin offen\hyp{}abgeschlossenen Mengen einen Basis für die Topologie $\tau$ bilden.
  \begin{itemize}
    \item[(1)] Ist $\mathcal{F}$ ein Filter ohne Adhärenzpunkt, so gilt für die Menge $\mathfrak{A} \subseteq \mathcal{F}$ aller offen\hyp{}abgeschlossenen Elemente von $\mathcal{F}$
     \begin{displaymath}
       \bigcap_{A \in \mathfrak{A}} A = \emptyset.
     \end{displaymath}
   \item[(2)] Erfüllt $(X,\tau)$ zusätzlich das T1\hyp{}Axiom, so existiert zu je zwei unterschiedlichen Punkten $x,y \in X$ ein $A_{(x,y)} \in \mathfrak{A}$ mit 
     \begin{displaymath}
       \{x,y\} \cap A_{(x,y)} = \{y\}.
     \end{displaymath}
  \end{itemize}
  %FUNKTIONIERT AUCH BEREITS MIT T0!!! S.a. Bartsch S.246
\end{lem}

\begin{proof}
  (1):
  Nach Voraussetzung besitzt $\mathcal{F}$ keinen Adhärenzpunkt.
  Gemäß Proposition \ref{prop:adherence} folgt $\bigcap_{F \in \mathcal{F}} \overline F = \emptyset$.
  Für alle $x \in X$ gilt damit, dass ein $F \in \mathcal{F}$ existiert, sodass $x \not\in \overline{F}$.
  Das bedeutet, dass $X \setminus \overline{F}$ eine offene Umgebung von $x$ ist.

  Da die offen\hyp{}abgeschlossenen Mengen nach Voraussetzung eine Basis für $\tau$ bilden, existiert eine offen\hyp{}abgeschlossene Umgebung $U$ von $x$, mit $U \subseteq X \setminus \overline{F}$ und damit $ U \cap F = \emptyset$.
  Folglich ist $X \setminus U \supseteq F$ und wegen des Abschlusses gegen Obermengenbildung $X \setminus U \in \mathcal{F}$.
  Es existiert also eine offen\hyp{}abgeschlossene Menge in $\mathcal{F}$, die $x$ nicht enthält. 
  Dies impliziert $ x \not\in \bigcap_{A \in \mathfrak{A}} A$.
  Da $x$ beliebig gewählt war, folgt $\bigcap_{A \in \mathfrak{A}} A = \emptyset$ und damit die Behauptung.

  (2):
  Seien nun $x,y \in X$ mit $x \neq y$ gegeben.
  Unter der Voraussetzung, dass $(X,\tau)$ das T1\hyp{}Axiom erfüllt, lässt sich zur offenen Umgebung $X \setminus \overline{F}$ aus Teil (1) des Beweises eine offene Umgebung 
  \begin{displaymath}
    U_x \subseteq X \setminus \overline{F}
  \end{displaymath}
  von $x$ finden, die zusätzlich $y$ nicht enthält. 
  Da die offen\hyp{}abgeschlossenen Mengen eine Basis von $\tau$ bilden, kann man zudem annehmen, dass $U_x$ offen\hyp{}abgeschlossen ist.
  Man prüft nach, dass für
  \begin{displaymath}
    A_{(x,y)} := X \setminus U_x
  \end{displaymath}
  die in (2) geforderten Eigenschaften erfüllt sind:
  Als Komplement einer offen\hyp{}abgeschlossenen Menge, ist auch $A_{(x,y)}$ offen und abgeschlossen.
  Nach Konstruktion ist 
  \begin{displaymath}
    x \not\in A_{(x,y)} \text{ sowie } y \in A_{(x,y)}.
  \end{displaymath}
  Zudem gilt, wie schon in (1) gezeigt, $A_{(x,y)} \in \mathcal{F}$.
  Daraus ergibt sich die Behauptung. \qedhere
\end{proof}

\begin{lem}
  \label{lem:openmap}
  Es seien $(X,\tau)$ und $(Y,\sigma)$ topologische Räume und es sei $\mathfrak{S}$ eine Subbasis für $\tau$.  
  Dann ist eine injektive Abbildung $f \colon X \to Y$ genau dann offen, wenn $f(\mathfrak{S}) \subseteq \sigma$.
\end{lem}

\begin{proof}
  Da $\mathfrak{S} \subseteq \tau$, folgt aus der Offenheit von $f$ sofort 
  \begin{displaymath}
    f(\mathfrak{S}) \subseteq f(\tau) \subseteq \sigma.
  \end{displaymath}

  Sei andererseits $f$ eine injektive Abbildung und $O \in \tau$ beliebig.
  Dann ist $O$ die Vereinigung endlicher Schnitte von Subbasislementen.
  Für zwei beliebige Mengen $O_1,O_2$ gilt nun 
  \begin{displaymath}
    f(O_1 \cap O_2) = f(O_1) \cap f(O_n)
  \end{displaymath}
  aufgrund der Injektivität von $f$, sowie 
  \begin{displaymath}
    f \left(\bigcup_{i \in I} O_i \right) = \bigcup_{i \in I} f(O_i)
  \end{displaymath}
  für eine beliebige Familie offener Mengen.
  Damit ist auch $f(O)$ offen und die Behauptung folgt.
\end{proof}

\begin{thm}
  \label{thm:ultrafiltercompact}
  Äquivalent sind:
  \begin{enumerate}
    \item[(1)] Der \textsc{Ascoli}-Satz bezüglich Ultrafilterkompaktheit.
    \item[(2)] \PIT.
  \end{enumerate}
\end{thm}

\begin{proof}
  $(1) \Rightarrow(2) $: Nach Satz \ref{thm:pitequivalence} reicht es aus zu zeigen, dass $(1)$ bereits impliziert, dass die \textsc{Kantor}\hyp{}Würfel $\mathbf{2}^I$ im gewöhnlichen Sinne kompakt sind.

  Angenommen, es existiere eine Menge $I$, sodass $P:= \mathbf{2}^I$ nicht kompakt ist.
  Man betrachte $I$ als einen mit der diskreten Topologie ausgestatteten topologischen Raum und $\mathbf{2}$ als diskreten metrischen Raum.
  Dann existiert nach Satz \ref{thm:compactness} ein Filter $\mathcal{F}$ auf $P$ ohne Adhärenzpunkt.

  Der mit der Produkttopologie ausgestattete Raum $P$ besitzt eine Basis aus offen\hyp{}abgeschlossenen Mengen.
  Dies folgt daraus, dass bereits jedes Subbasiselement $S$ per Definition offen ist und, da $\mathbf{2}$ diskret ist, auch ein offenes Komplement besitzt.
  Damit ist $S$ aber als Komplement einer offenen Menge definitionsgemäß abgeschlossen.
  Insgesamt ist $S$ also offen\hyp{}abgeschlossen.
  Da die Menge der offen\hyp{}abgeschlossenen Mengen die endliche Durchschnittseigenschaft besitzt, ist auch jedes von der Subbasis erzeugte Basiselement eine offen\hyp{}abgeschlossene Menge.

  Hiermit lässt sich nun Lemma \ref{lem:filterclusterpoint}(1) auf $\mathcal{F}$ anwenden.
  Damit folgt für die Menge $\mathfrak{A}$ aller offen\hyp{}abgeschlossenen Elemente von $\mathcal{F}$
  \begin{displaymath}
    \bigcap_{A \in \mathfrak{A}} A = \emptyset \tag{$\star$}.
  \end{displaymath}

  Man kann nun für alle $A \in \mathfrak{A}$ die charakteristische Funktion
  \begin{equation}
    \label{eq:clopenmap}
    f_A \colon P \to \mathbf{2}, \quad
    f_A(x) := 
    \begin{cases}
      1 &\text{, falls } x \in A \\
      0 &\text{, falls } x \not\in A
    \end{cases}
  \end{equation}
  definieren.
  Da alle $A$ offen\hyp{}abgeschlossen sind, folgt sofort die Stetigkeit von $f_A$ für alle $A \in \mathfrak{A}$.
  Somit induziert die Familie $(f_A)_{A \in \mathfrak{A}}$ eine Abbildung 
  \begin{displaymath}
    f \colon P \to \mathbf{2}^\mathfrak{A}, \quad
    x \mapsto (f_A(x))_{A \in \mathfrak{A}},
  \end{displaymath}
  welche bezüglich der Produkttopologie auf $\mathbf{2}^\mathfrak{A}$ auf $P$ stetig ist.

  Es bezeichne 
  \begin{displaymath}
    F := f(P)
  \end{displaymath}
  das Bild von $P$ unter obiger Abbildung $f$.
  Es soll nun gezeigt werden, dass $f$ eine Einbettung ist, also ein Homöomorphismus auf $F$:
  Die Abbildung $f$ ist injektiv, denn für $x,y \in P$ mit $x \neq y$ gilt, da $P$ als \textsc{Hausdorff}\hyp{}Raum insbesondere das T1\hyp{}Axiom erfüllt,
  \begin{displaymath}
    f_{A_{(x,y)}}(x) = 0 \neq 1 = f_{A_{(x,y)}}(y),
  \end{displaymath}
  mit $A_{(x,y)} \in \mathfrak{A}$ aus Lemma \ref{lem:filterclusterpoint}(2).

  Zudem ist $f$ offen bezüglich der Spurtopologie auf $F$, wie sich mit Lemma \ref{lem:openmap} beweisen lässt. 
  Dazu sei nun $O = \pi_{i_0}^{-1}(\{j\})$ mit $i_0 \in I$ und $j \in \{0,1\}$ eine beliebige offene Subbasismenge der Produkttopologie auf $P$,  wobei die kanonischen Projektionen des Produktraumes mit $\pi_i$ bezeichnet seien.
  Man erkennt $O$ als eine offen\hyp{}abgeschlossene Menge aufgrund der Stetigkeit der kanonischen Projektionen.

  Gilt $O \in \mathcal{F}$, also insbesondere $O \in \mathfrak{A}$, so ist für alle $x \in O$
  \begin{displaymath}
    \pi_O(f(x)) = f_O(x) = 1.
  \end{displaymath}
  Daraus folgt 
  \begin{displaymath}
    f(O) = \pi_O^{-1}(\{1\}) \cap F
  \end{displaymath}
  und folglich ist $f(O)$ offen bezüglich der Teilraumtopologie auf $F$.

  Ist andernfalls $O \not\in \mathcal{F}$, so gilt für alle $M \in \mathcal{F}$ entweder
  \begin{displaymath}
    M \cap O \neq \emptyset
  \end{displaymath}
  oder $P \setminus O \in \mathcal{F}$.
  Im zweiten Fall gilt damit insbesondere $P \setminus O \in \mathfrak{A}$, da $O$ offen\hyp{}abgeschlossen ist, und analog zum ersten Teil folgt für alle $x \in O$
  \begin{displaymath}
    \pi_{P \setminus O}(f(x)) = 0
  \end{displaymath}
  und damit 
  \begin{displaymath}
    f(O) = \pi_{P \setminus O}^{-1}(\{0\}) \cap F.
  \end{displaymath}
  Im anderen Fall gilt zudem für alle $M \in \mathcal{F}$
  \begin{displaymath}
    M \cap (P \setminus O) \neq \emptyset
  \end{displaymath}
  und damit 
  \begin{displaymath}
    f(O) = F,
  \end{displaymath}
  denn in $O$ existiert zu jedem $A \in \mathfrak{A}$ ein $x_1 \in O$ mit $f_A(x) = 1$ sowie ein $x_2 \in O$ mit $f_A(x) = 0$.
  Also ist $f(O)$ auch in diesem Fall offen bezüglich Spurtopologie.  
  Damit ist Lemma \ref{lem:openmap} anwendbar und impliziert, dass $f$ eine offene Abbildung auf ihr Bild ist.
  Folglich ist $f$ ein Homöomorphismus auf $F$.

  Jeder endliche diskrete Raum ist ultrafilterkompakt, da dort die einzigen Ultrafilter, die Einpunktfilter sind, welche immer konvergieren.
  Es ist daher $P$ nach Satz \ref{lem:uftprod} als Produkt ultrafilterkompakter Räume wiederum ultrafilterkompakt.
  Es folgt, dass $F$ als homöomorphes Bild eines ultrafilterkompakten Raumes ebenso ultrafilterkompakt ist.

  Es soll nun den \textsc{Ascoli}\hyp{}Satz bezüglich Ultrafilterkompaktheit auf $F$ angewendet werden.
  Betrachtet man $\mathfrak{A}$ als diskreten topologischen Raum und $\mathbf{2}$ als diskreten metrischer Raum, so ist Bedingung (a) von Satz \ref{thm:ascoliPIT} erfüllt, denn, da in diskreten Räumen eine Teilmenge genau dann kompakt ist, wenn sie endlich ist, folgt, dass die Subbasis der kompakt\hyp{}offenen Topologie auf $\mathbf{2}^\mathfrak{A}$ genau mit der Subbasis der punktweisen Topologie übereinstimmt. 
  Andererseits ist die Aussage $(b, \beta)$ nicht erfüllt:
  Es ist zwar $p := (1)_{A \in \mathfrak{A}} \in \mathbf{2}^\mathfrak{A}$ im Abschluss von $F$ in $\mathbf{2}^\mathfrak{A}$ enthalten, denn für eine beliebige Umgebung $U_p \subseteq \mathbf{2}^\mathfrak{A}$ von $p$, existieren Subbasiselemente $B_{A_1},\dots,B_{A_n}$ mit 
  \begin{displaymath}
    p \in \bigcap_{i=1}^n B_{A_i} \subseteq U_p,
  \end{displaymath}
  wobei $B_{A_i} = \prod_{A \in \mathfrak{A}} R_{A}$ mit  $R_{A_i} = \{1\}$ und $R_A = \mathbf{2}$ sonst gelte.
  Dies wiederum impliziert
  \begin{displaymath}
    \bigcap_{i=1}^n B_{A_i} \cap F \neq \emptyset,
  \end{displaymath}
  denn für 
  \begin{displaymath}
    x \in \bigcap_{i=1}^n A_i \in \mathcal{F}
  \end{displaymath}
  gilt $f_{A_i}(x) = 1$, weil letzlich offen\hyp{}abgeschlossene Mengen sowie Filter die endliche Durchschnittseigenschaft besitzen.
  Also gilt auch $U_p \cap F \neq \emptyset$ und damit $p \in \overline F$.
  
  Andererseits ist $p$ jedoch nicht in $F$ enthalten. 
  Gälte nämlich $p \in F$, so existiert ein $x \in P$ mit $f_A(x) = 1$ für alle $A \in \mathfrak{A}$, was gleichbedeutend ist mit $x \in \bigcap_{A \in \mathfrak{A}} A$.
  Dies steht jedoch im Widerspruch dazu, dass nach ($\star$) der Schnitt über alle in $\mathcal{F}$ enthaltenen offen\hyp{}abgeschlossenen Mengen leer ist.
  Also gilt $F \neq \overline{F}$, ein Widerspruch zum \textsc{Ascoli}\hyp{}Satz bezüglich Ultrafilterkompaktheit.
  Die Annahme, dass $\mathbf{2}^I$ nicht kompakt ist, muss also verworfen werden.

  (2) $\Rightarrow$ (1): \PIT impliziert nach Satz \ref{thm:uftcompact} die Übereinstimmung des gewöhnlichen Kompaktheitsbegriffes mit dem Begriff der Ultrafilterkompaktheit. Somit folgt (1) direkt aus Satz \ref{thm:ascoliPIT}.
\end{proof}

Ähnliches gilt nun auch für den nächsten Kompaktheitsbegriff der \textsc{Tychonoff}\hyp{}Kompaktheit wie der folgende Satz zeigt.

\begin{thm}
  Äquivalent sind:
  \begin{enumerate}
    \item[(1)] Der \textsc{Ascoli}-Satz bezüglich \textsc{\textsc{Tychonoff}}\hyp{}Kompaktheit.
    \item[(2)] \PIT.
  \end{enumerate}
\end{thm}

\begin{proof}
  (1) $\Rightarrow$ (2): Es reicht aus zu zeigen, dass (1) bereits impliziert, dass die \textsc{Kantor}\hyp{}Würfel $\mathbf{2}^I$ im gewöhnlichen Sinne kompakt sind.
  Angenommen, es existiert eine Menge $I$, sodass $P:= \mathbf{2}^I$ nicht kompakt ist.
  Man betrachte $I$ als einen mit der diskreten Topologie ausgestatteten topologischen Raum und $\mathbf{2}$ als diskreten metrischen Raum.
  Dann existiert ein Filter $\mathcal{F}$ auf $P$ ohne Adhärenzpunkt.

  Da das Intervall $[0,1]$ ausgestattet mit euklidischer Topologie insbesondere ein T1\hyp{}Raum ist, sind Einpunktmengen abgeschlossen.
  Betrachtet man nun $P$ als Teilraum des \textsc{Hilbert}\hyp{}Würfels $[0,1]^I$, so folgt nach Proposition \ref{prop:cartesianclosed}, dass $P$ abgeschlossen im \textsc{Hilbert}\hyp{}Würfel ist.
  Definitionsgemäß ist also $P$ \textsc{Tychonoff}\hyp{}kompakt.

  Wie aus dem Beweis von Satz \ref{thm:ultrafiltercompact} folgt, ist $P$ homöomorph zu $F:= f(P)$, wobei $f$ die auf der Menge der offen\hyp{}abgeschlossenen Elemente von $\mathcal{F}$ definierte Funktion aus Gleichung (\ref{eq:clopenmap}) sei.
  Somit ist $F$ ebenfalls \textsc{Tychonoff}\hyp{}kompakt.

  Wendet man nun den \textsc{Ascoli}\hyp{}Satz bezüglich \textsc{Tychonoff}\hyp{}Kompaktheit auf $F$ an, so folgt anlog zu Satz \ref{thm:ultrafiltercompact}, dass zwar die Bedingung (a) gilt aber die Aussage (b, $\beta$) nicht erfüllt ist.

  (2) $\Rightarrow$ (1): \PIT impliziert die Übereinstimmung des gewöhnlichen Kompaktheitsbegriffes mit dem Begriff der \textsc{Tychonoff}\hyp{}Kompaktheit. Somit folgt (1) direkt aus Satz \ref{thm:ascoliPIT}.
\end{proof}

Die Substitution der Kompaktheit durch \textsc{Alexandroff}\hyp{}\textsc{Urysohn}\hyp{}Kompaktheit führt sogar zu einer Verschärfung der benötigten Bedingungen für die Gültigkeit des \textsc{Ascoli}\hyp{}Satzes wie der nächste Satz zeigt.

\begin{thm}
  Äquivalent sind:
  \begin{enumerate}
    \item[(1)] Der \textsc{Ascoli}-Satz bezüglich \textsc{Alexandroff}\hyp{}\textsc{Urysohn}\hyp{}Kompaktheit.
    \item[(2)] $\mathbf{AC}$.
  \end{enumerate}
\end{thm}

\begin{proof}
  (1) $\Rightarrow$ (2): Sei $X$ eine Menge und $\mathbf{2}$ die zweielementige Menge $\{0,1\}$ und seien beide mit der diskreten Topologie ausgestattet. 
  Die diskrete Topologie lässt sich auch als die von der \textit{diskreten Metrik} induzierte Topologie auffassen, somit kann $\mathbf{2}$ auch als metrischer Raum betrachtet werden. 

  Es sei $F := C(X, \mathbf{2})$.
  Dann gilt $C(X, \mathbf{2}) = \mathbf{2}^X$, da $X$ nach Voraussetzung diskret ist.
  Da $\mathbf{2} = F(x)$ für alle $x \in X$ als endlicher Raum keine unendliche Teilmenge besitzt, folgt sofort die \textsc{Alexandroff}\hyp{}\textsc{Urysohn}\hyp{}Kompaktheit von $F(x)$.
  Nach Definition ist $\mathbf{2}^X$ als Gesamtraum abgeschlossen bezüglich jeder Topologie, also insbesondere auch bezüglich der punktweisen Topologie $\tau_p$.
  Zuletzt ist $F$ auch gleichgradig stetig, da $X$ diskret und somit Einpunktmengen offen sind, d.h. für alle $x \in X$ und $\varepsilon > 0$ wähle man $\{x\}$ als offene Umgebung in $X$. 
  Dann gilt für alle $f \in F$, aufgrund der Definitheit der Metrik für  $d(f(x),f(x)) = 0 \leq \varepsilon$.
  Damit sind die Bedingungen (b) des \textsc{Ascoli}\hyp{}Satzes erfüllt und es folgt, dass $\mathbf{2}^X$ \textsc{Alexandroff}\hyp{}\textsc{Urysohn}\hyp{}kompakt ist. 

  Es soll Satz \ref{thm:comparablecardinal} angewendet werden, um zu folgern, dass \AC gilt.
  Dazu seien $A$ und $B$ zwei unendliche Mengen mit Kardinalitäten $|A| = a$ und $|B| = b$.
  Sei $\infty$ ein Element, welches nicht in $A \cup B$ enthalten ist.
  Setze $I := \{\infty\} \cup A \cup B$ und betrachte den \textsc{Kantor}\hyp{}Würfel $\mathbf{2}^I$.
  Nach dem vorangehenden Beweis ist dieser \textsc{Alexandroff}\hyp{}\textsc{Urysohn}\hyp{}kompakt.
  Für jede Teilmenge $C \subseteq I$ sei nun 
  \begin{equation}
    \chi_C \colon I \to \mathbf{2}, \quad
    \chi_C(i) := 
    \begin{cases}
      1 &\text{, falls } i \in C \\
      0 &\text{, falls } i \not\in C
    \end{cases}
  \end{equation}
  die zu $C$ gehörige charakteristische Funktion.
  Dan besitzt die Teilmenge
  \begin{displaymath}
    M := \{\chi_{\{a\}} \mid a \in A\} \cup \{\chi_{\{\infty,b\}} \mid b \in B\} \subseteq \mathbf{2}^I
  \end{displaymath}
  als unendliche Menge einen vollständigen Häufungspunkt $x := (x_i)_{i \in I} \in \mathbf{2}^I$.

  Man unterscheide nun zwei Fälle:
  Ist $x_\infty = 0$, so ist das Subbasiselement $U := \prod_{i \in I} U_i$ mit $U_\infty = \{0\}$ und $U_i = \mathbf{2}$ für alle $i \neq \infty$ eine offene Umgebung von $x$. 
  Des Weiteren gilt
  \begin{displaymath}
    |A|
    = |\{\chi_{\{a\}} \mid a \in A\}|
    = | M \cap U |
    = | M |
  \end{displaymath}
  da $x$ ein vollständiger Häufungspunkt von $M$ ist.
  Es existiert also eine Bijektion von $A$ nach $M$, insbesondere existiert eine Injektion von $\{\chi_{\{\infty,b\}} \mid b \in B\}$ nach $A$, und damit gilt definitionsgemäß $b \leq a$. 
  Gilt andererseits $x_\infty = 1$, so zeigt man analog $a \leq b$.
  Da $a$ und $b$ beliebige unendliche Kardinalitäten waren und endliche Kardinalitäten immer vergleichbar sind, findet Satz \ref{thm:comparablecardinal} Anwendung.

  (2) $\Rightarrow$ (1): Das Auswahlaxiom impliziert nach Satz \ref{thm:alexandroffUrysohnCompactAC} die Übereinstimmung des gewöhnlichen Kompaktheitsbegriffes mit dem Begriff der \textsc{Alexandroff}\hyp{}\textsc{Urysohn}\hyp{}Kompaktheit. Somit folgt (1) direkt aus Satz \ref{thm:ascoliPIT}.
\end{proof}

\section{Der klassische \textsc{Ascoli}\hyp{}Satz und eine Modifikation}

Wie die vorangehenden Ausführungen gezeigt haben, versagt der \textsc{Ascoli}\hyp{}Satz bei allen vorgestellten Versionen von Kompaktheit, falls lediglich \ZF vorausgesetzt wird. 
Es soll im Folgenden untersucht werden, inwiefern eine eingeschränkte Formulierung des \textsc{Ascoli}\hyp{}Satzes die Anforderungen senkt.

\begin{defn}[Klassischer \textsc{Ascoli}-Satz] 
  \label{def:klassisch}
  Für eine Funktionenmenge $F$ stetiger Abbildungen $f \colon \R \to \R$ sind die folgenden Bedingungen äquivalent:
  \begin{itemize}
    \item[(a)] Jede Folge $(f_n)_{n \in \N}$ in $F$ besitzt eine Teilfolge $(f_{(\nu(n))})_{n \in \N}$, die \textit{stetig gegen} eine nicht notwendig in $F$ liegende Funktion $g$ \textit{konvergiert}. Das bedeutet:
      \begin{addmargin}[2em]{2em}% 1em left, 2em right
        Für alle $x \in \R$ und alle $(x_n)_{n \in \N} \in \R^\N$ gilt: \\
        Aus $\lim_{n \to \infty} x_n = x$ folgt $\lim_{n \to \infty} f_{\nu(n)}(x_n) = g(x)$.
      \end{addmargin}
    \item[(b)]
      \begin{itemize}
        \item[($\alpha$)] Für alle $x \in \R$ ist die Menge $F(x) = \{ f(x) \mid f \in F\}$ beschränkt.
        \item[($\beta$)] Es ist $F$ gleichgradig stetig.
      \end{itemize}
  \end{itemize}
\end{defn}

Dass die in Definition \ref{def:klassisch} auftretende Grenzfunktion $g$ eindeutig bestimmt ist, sieht man über ein Folgenmischungsargument ein:
Sind nämlich $(x_n)_{n \in \N}$ und $(y_n)_{n \in \N}$ zwei Folgen in $\R$ mit Grenzwert $x$, so lässt sich eine dritte Folge $(z_n)_{n \in \N}$ konstruieren mit $z_{2n-1} := x_{2n-1}$ und $z_{2n} := y_{2n}$.
Diese konvergiert weiterhin gegen $x$ und es gilt 
\begin{displaymath}
   \lim_{n \to \infty} f_{n}(x_n) 
  =\lim_{n \to \infty} f_{2n-1}(x_{2n-1}) 
  =\lim_{n \to \infty} f_{2n-1}(z_{2n-1}) 
  =\lim_{n \to \infty} f_n(z_n),
\end{displaymath}
da jede Teilfolge einer konvergenten Folge gegen denselben Grenzwert konvergiert.
Für $(y_n)_{n \in \N}$ gilt analoge Gleichung, was letztlich
\begin{displaymath}
   \lim_{n \to \infty} f_{n}(x_n) = \lim_{n \to \infty} f_{n}(y_n) 
\end{displaymath}
impliziert.

Es sollen zunächst einige Eigenschaften von Funktionenfolgen in Bezug auf die stetige Konvergenz festgehalten werden.

\begin{prop}
  \label{prop:stetigkonv}
  Es sei $(f_n)_{n \in \N}$ eine Folge stetiger Abbildungen von $\R$ nach $\R$ und $g$ eine stetige Abbildung mit demselben Definitionsbereich und derselben Zielmenge.
  Dann gelten folgende Aussagen:
  \begin{itemize}
    \item[(1)] Falls $(f_n)_{n \in \N}$ stetig gegen $f$ konvergiert, so konvergiert $(f_n)_{n \in \N}$ auch punktweise gegen $f$.
    \item[(2)] Falls $(f_n)_{n \in \N}$ stetig gegen $f$ konvergiert, so ist auch $f$ stetig.
    \item[(3)] Falls $(f_n)_{n \in \N}$ lokal gleichmäßig gegen $f$ konvergiert, so konvergiert $(f_n)_{n \in \N}$ auch stetig gegen $f$.
  \end{itemize}
\end{prop}

\begin{proof}
  (1): 
  Dies folgt direkt aus der Definition unter Betrachtung konstanter Folgen $(x_n)_{n \in \N}$ mit $x_n = x$ für alle $n \in \N$.

  (2): 
  Es soll die Stetigkeit der Grenzfunktion $f$ mittels Folgenstetigkeit gezeigt werden. 
  Es sei dazu $(x_n)_{n \in \N}$ eine Folge in $\R$ mit Grenzwert $x$ und $\varepsilon > 0$ gegeben.
  Nach (1) konvergiert $(f_n)_{n \in \mathbb{N}}$ punktweise.
  Daraus lässt sich nun induktiv eine Teilfolge $(f_{n_k})_{k \in \mathbb{N}}$ gewinnen mit
  \begin{displaymath}
    | f_{n_k}(x_k) - f(x_k) | \leq \frac{\varepsilon}{2}, \text{ für alle $k \in \N$},
  \end{displaymath}
  denn aufgrund der punktweisen Konvergenz existiert zu $x_1$ ein $n_1$, sodass
  \begin{displaymath}
    | f_{n}(x_1) - f(x_1) | \leq \frac{\varepsilon}{2}, \text{ für alle $n \geq n_1$}.
  \end{displaymath}
  Analog findet man für $x_k$ ein $n_k > n_{k-1}$ mit einer analogen Eigenschaft.

  Die so konstruierte Teilfolge besitzt nach wie vor denselben Grenzwert
  \begin{displaymath}
    \lim_{k \to \infty} f_{n_k}(x_k) = f(x).
  \end{displaymath}
  Man sieht dies ein, indem man eine neue Folge $(y_m)_{m \in \N}$ definiert über $y_m:=x_1$ für $1 \leq m \leq n_1$ und $y_m:=x_k$ für $n_{k-1} < m \leq n_{k}$ für $k > 1$, sodass $\lim_{m \to \infty} y_m = x$.
  Damit gilt nun
  \begin{displaymath}
    \lim_{m \to \infty} f_m(y_m) = f(x)
  \end{displaymath}
  aufgrund der stetigen Konvergenz der Folge $(f_n)_{n \in \N}$ und somit
  \begin{displaymath}
    \lim_{k \to \infty} f_{n_k}(x_{k}) 
    = \lim_{k \to \infty} f_{n_k}(y_{n_k}) 
    = f(x).
  \end{displaymath} 

  Es existiert also ein $n_\varepsilon \in \N$, sodass für alle $n_k \geq n_\varepsilon$ gilt
  \begin{displaymath}
    | f_{n_k}(x_k) - f(x) | \leq \frac{\varepsilon}{2}.
  \end{displaymath}
  Daraus folgt nun unter Anwendung der Dreiecksungleichung für alle $n \geq n_\varepsilon$
  \begin{displaymath}
    | f(x_k) - f(x) |
    \leq
    | f_{n_k}(x_k) - f(x_k) | + | f_{n_k}(x_k) - f(x) | 
    \leq \varepsilon.
  \end{displaymath}

  (3): 
  Es sei $(x_n)_{n \in \N}$ eine Folge in $\R$ mit Grenzwert $x$.
  Des Weiteren sei $\varepsilon > 0$ gegeben.
  Da nach Voraussetzung $(f_n)_{n \in \N}$ lokal gleichmäßig gegen $f$ konvergiert existiert eine offene Umgebung $U_x$ von $x$, sodass ${f_n}_{| U_x}$ gleichmäßig konvergiert.
  Es existiert somit ein $n_\varepsilon$, sodass für alle $y \in U_x$ gilt 
  \begin{displaymath}
    | f_n(y) - f(y) | \leq \frac{\varepsilon}{2} \quad \text{ für alle $n \geq n_\varepsilon$}.
  \end{displaymath}
  Als Grenzfunktion einer lokal gleichmäßig konvergenten Funktionenfolge ist $f$ zudem stetig in $x$.
  Es existiert also ein $\delta > 0$, sodass für alle $y \in U_\delta(x)$
  \begin{displaymath}
    | g(y) - g(x) | < \frac{\varepsilon}{2}
  \end{displaymath}
  gilt.
  Da $(x_n)_{n \in \N}$ voraussetzungsgemäß gegen $x$ konvergiert, existiert ein $N \in \N$, sodass $x_n \in U_\delta(x) \cap U_x$ für alle $n \geq N$.

  Setzt man nun $n':= \max(n_\varepsilon, N)$, so folgt unter Anwendung der Dreiecksungleichung
  \begin{displaymath}
    | f_n(x_n) - f(x) |
    \leq
    | f_n(x_n) - f(x_n) | + | f(x_n) - f(x) | 
    \leq \varepsilon \text{ für alle $n \geq n'$}. \qedhere
  \end{displaymath}
\end{proof}

Für die anschließende mengentheoretische Charakterisierung des klassischen \textsc{Ascoli}\hyp{}Satzes wird folgender Satz benötigt.

\begin{thm}
  \label{thm:unboundedsubsetofR}
  Folgende Aussagen sind äquivalent:
  \begin{enumerate}[(1)]
    \item \CCR, das Axiom der abzählbaren Auswahl (Countable Choice) für Teilmengen reeller Zahlen:
      \begin{addmargin}[2em]{2em}% 1em left, 2em right
        \textit{Für jede Folge $(X_n)_{n \in \N}$ nichtleerer Teilmengen $X_n \subseteq \R$ ist das kartesische Produkt $\prod_{n \in \N} X_n$ nichtleer.}
      \end{addmargin}
    \item \PCCR, das Axiom der partiellen abzählbaren Auswahl (Partial Countable Choice) für Teilmengen reeller Zahlen: 
      \begin{addmargin}[2em]{2em}% 1em left, 2em right
        \textit{Für jede Folge $(X_n)_{n \in \N}$ nichtleerer Teilmengen $X_n \subseteq \R$ existiert eine unendliche Teilmenge $M \subseteq \N$, sodass das kartesische Produkt $\prod_{m \in M} X_m$ nichtleer ist.}
      \end{addmargin}
    \item Jede unbeschränkte Teilmenge der reellen Zahlen $\R$ enthält eine unbeschränkte Folge.
  \end{enumerate}
\end{thm}

\begin{proof}
  (1) $\Leftrightarrow$ (2):
  Es ist klar, dass aus \CCR die Aussage \PCCR folgt.
  Sei umgekehrt eine Folge $(X_n)_{n \in \N}$ nichtleerer Teilmengen der reellen Zahlen gegeben.
  Man definiere nun $Y_n := \prod_{m \leq n} X_m$.
  Da endliche kartesische Produkte nichtleerer Mengen wiederum nichtleer sind, ist $(Y_n)_{n \in \N}$ eine Folge nichtleerer Mengen.
  Nach \PCCR existiert eine unendliche Teilmenge $M \subseteq \N$ und ein Element $(y_m)_{m \in \N} \in \prod_{m \in M} Y_m$.
  Daher gilt $y_m \in Y_m = \prod_{k \leq m} X_k$, also $y_m = (x_{1,m},\dots,x_{m,m})$.

  Man definiere nun für alle $n \in \N$ die Funktion
  \begin{displaymath}
    m(n) := \min\{m \in M \mid n \leq m\}.
  \end{displaymath}
  Sie ist aufgrund der Wohlordnung der natürlichen Zahlen wohldefiniert.
  Dann ist das Element $(x_{n,m(n)})_{n \in \N} \in \prod_{n \in \N} X_n$.
  Daraus folgt die Behauptung.

  (1) $\Rightarrow$ (3):
  Sei $X \subseteq \R$ eine nichtleere unbeschränkte Teilmenge der reellen Zahlen.
  Man betrachte die Mengenfolge $X_n := \{x \in X \mid x \geq n, n \in \N\}$.
  Nach Voraussetzung gilt $X_n \neq \emptyset$ für alle $n \in \N$.
  Die Anwendung von \CCR ergibt folglich ein $(x_n)_{n \in \N} \in \prod_{n \in \N} X_n$.
  Nach Konstruktion ist $(x_n)_{n \in \N}$ eine unbeschränkte Folge in $X$.

  (3) $\Rightarrow$ (2):
  Sei $(X_n)_{n \in \N}$ eine Folge nichtleerer Teilmengen reeller Zahlen.
  Sei zusätzlich $f \colon \R \to (0,1)$ eine fest gewählte Bijektion.
  Ein Beispiel hierfür liefert die Funktion $f(x) = \frac{2}{\pi}\arctan(x)$.
  Des Weiteren definiere man für alle $n \in \N$
  \begin{displaymath}
    \sigma_n \colon \R \to \R, \quad \sigma_n(x) := n + x.
  \end{displaymath}
  Diese Abbildung ist injektiv.
  Man setze nun $Y_n := \sigma_n(f(X_n))$. 
  Dann folgt aus $X_n \subseteq \R$ sofort $Y_n \subseteq (n,n+1)$ und die Menge $Y:= \bigcup_{n \in \N} Y_n$ ist eine unbeschränkte nichtleere Teilmenge von $\R$, da aus $X_n \neq \emptyset$ stets $Y_n \neq \emptyset$ folgt.
  Nach Voraussetzung existiert eine unbeschränkte Folge $(y_n)_{n \in \N}$ in $Y$.
  Damit definiert
  \begin{displaymath}
    M := \{m \in \N \mid \text{Es existiert ein } n \in \N \text{ mit } y_n \in Y_m\}
  \end{displaymath}
  Es ist $M$ eine unendliche Teilmenge natürlicher Zahlen.
  Angenommen, $M$ wäre endlich.
  Dann würde ein $m = \max(M)$ existieren, sodass $m + 1$ eine obere Schranke für die Folge $(y_n)_{n \in \N}$ ist, was jedoch der Unbeschränktheit der Folge widerspricht.
  Zudem ist $\N$ mit der natürlichen Ordnungsrelation wohlgeordnet.
  Daraus folgt $Y := \prod_{m \in M} Y_m \neq \emptyset$, indem man $g(m) := \min\{n \in \N \mid y_n \in Y_m\}$ setzt.
  Es bezeichne $(\tilde{y}_m)_{m \in M}$ ein Element aus $Y$.
  Aufgrund der Injektivität von $f$ und $\sigma_m$ existiert ein eindeutig bestimmtes Element $x_m \in X_m$ mit $\sigma_m(f(x_m)) = \tilde{y}_m$.
  Dann ist $(x_m)_{m \in M} \in \prod_{m \in M} X_m$.
  Dies ist jedoch gerade die Aussage von \PCCR.
\end{proof}

\begin{bem}
  Man erkennt, dass im Beweis von Satz \ref{thm:unboundedsubsetofR}((1) $\Leftrightarrow$ (2)) keine Eigenschaften der reellen Zahlen verwendet wurden. Der Äquivalenzbeweis von \CC und \PCC für beliebige abzählbare Mengenfamilien erfolgt daher analog.
\end{bem}

\begin{thm}
  \label{thm:classicalascoli}
  Äquivalent sind:
  \begin{enumerate}
    \item[(1)] Der klassische \textsc{Ascoli}-Satz.
    \item[(2)] \CCR.
  \end{enumerate}
\end{thm}

\begin{proof}
  (1) $\Rightarrow$ (2):
  Nach Satz \ref{thm:unboundedsubsetofR} genügt es zu zeigen, dass jede unbeschränkte Teilmenge $B$ von $\R$ eine unbeschränkte Folge enthält.
  Dazu sei $B \subseteq \R$ unbeschränkt.
  Für ein $b \in B$ lassen sich die konstante Abbildung $f_b \colon \R \to \R, f(x):= b$ und die Menge $F := \{f_b \mid b \in B\}$ definieren.
  Für $x \in \R$ gilt $F(x) = \{f_b(x) \mid b \in B\} = B$ und diese Menge ist nach Voraussetzung unbeschränkt.
  Also ist Teil (b, $\alpha$) des klassischen \textsc{Ascoli}\hyp{}Satzes verletzt.
  Mit (1) folgt sogleich, dass auch Teil (a) des klassischen \textsc{Ascoli}\hyp{}Satzes nicht gelten kann und somit die Existenz einer Folge $(f_{b_n})_{n \in \N}$ mit $f_{b_n} = b_n$ ohne stetig konvergente Teilfolge.
  Daraus folgt, dass die Folge $(b_n)_{n \in \N}$ unbeschränkt ist.

  (2) $\Rightarrow$ (1):
  Es sei $F$ eine Menge stetiger Abbildungen $f \colon \R \to \R$.

  (a) $\Rightarrow$ (b,$\alpha$):
  Angenommen (b,$\alpha$) gelte nicht. So existiert ein $x \in \R$, sodass $F(x)$ unbeschränkt ist.
  Nach Definition existiert also für alle $C \in \R$ ein $f \in F$ mit $|f(x)| \geq C$.
  Insbesondere sind also die Mengen 
  \begin{displaymath}
    F_n := \{ f \in F \mid |f(x)| \geq m \}
  \end{displaymath}
  für alle $n \in \N$ nichtleer.
  \CCR impliziert folglich die Existenz einer Folge 
  \begin{displaymath}
    (f_n)_{n \in \N} \in \prod_{n \in \N} F_n.
  \end{displaymath}
  Für obiges $x \in \R$ gilt also insbesondere $|f_{n_k}(x)| \geq n_k \geq n$ für alle $k \in \N$, sodass die Funktionenfolge $(f_n)_{n \in \N}$ in $x$ nicht punktweise, folglich auch nicht stetig konvergiert.
  Dies widerspricht jedoch Bedingung (a).

  (a) $\Rightarrow$ (b,$\beta$):
  Angenommen (b, $\beta$) gelte nicht.
  So existiert ein $x \in \R$ und $\varepsilon > 0$, sodass das Bild keiner offenen Umgebung von $x$ in $U_\varepsilon(f(x))$ enthalten ist.
  Insbesondere gilt also auch, dass für kein $\delta > 0$ gilt, dass
  \begin{displaymath}
    f(U_\delta(x)) \subseteq U_\varepsilon(f(x)).
  \end{displaymath}
  Schließlich gilt für alle $n \in \N$, dass die Mengen
  \begin{displaymath}
    F_n := \{ f \in F \mid \text{Es existiert ein } y \in U_{\frac{1}{n+1}}(x) \text{ mit } f(y) \not\in U_\varepsilon(f(x)) \}
  \end{displaymath}
  nichtleer sind.
  Wie zuvor impliziert \CCR die Existenz einer Folge 
  \begin{displaymath}
    (f_n)_{n \in \N} \in \prod_{n \in \N} F_n.
  \end{displaymath}
  Nimmt man an, dass diese Folge eine stetig konvergente Teilfolge $(f_{\nu(n)})_{n \in \N}$ besitzt, so konvergiert diese nach Proposition \ref{prop:stetigkonv}(a) auch punktweise.
  Zu dem anfangs gegebenen $x$ und $\varepsilon > 0$ existiert also insbesondere ein $N \in \N$, sodass
  \begin{displaymath}
    | f_{\nu(n)}(x) - g(x) | \leq \frac{\varepsilon}{2}.
  \end{displaymath}
  Andererseits existiert nach der Definition von $F_n$ eine Folge $(y_n)_{n \in \N}$ mit Grenzwert $x$ und
  \begin{displaymath}
    | f_{\nu(n)}(x) -  f_{\nu(n)}(y_n) | \geq \varepsilon.
  \end{displaymath}
  Unter Verwendung der inversen Dreiecksungleichung folgt damit für alle $n \geq N$
  \begin{displaymath}
    | f_{\nu(n)}(y_n) - g(x) | 
    \geq 
    | f_{\nu(n)}(y_n) - f_{\nu(n)}(x) | - | f_{\nu(n)}(x) - g(x)  | 
    \geq \frac{\varepsilon}{2}.
  \end{displaymath}
  Dies steht jedoch im Widerspruch zur stetigen Konvergenz von $(f_{\nu(n)})_{n \in \N}$.

  (b) $\Rightarrow$ (a):
  Sei $(f_n)_{n \in \N}$ eine Folge in $F$.
  Des Weiteren sei $(r_n)_{n \in \N}$ eine Abzählung der rationalen Zahlen.
  Im Folgenden soll induktiv eine Folge von geordneten Paaren $(a_n, s_n)_{n \in \N}$ mit $a_n \in \R$ für alle $n \in \N$ und eine Folge $s_n = (g_m^n)_{m \in \N}$ in $F$ definiert werden:
  \begin{enumerate}
    \item 
      Nach $(b)$ ist $F(x)$ für alle $x \in \R$ beschränkt. 
      Somit ist auch 
      \begin{displaymath}
        \{f_n(r_0) \mid n \in \N\} \subseteq F(r_0)
      \end{displaymath}
      beschränkt.
      Man setze nun $a_0 := \liminf_{n \to \infty} f_n(r_0)$.
      Weiterhin definiere man $s_0 = (g_n^0)_{n \in \N}$ induktiv als Teilfolge $(f_{\nu(n)})_{n \in \N}$ von $(f_n)_{n \in \N}$:
      \begin{itemize}

        \item[$a)$] $\nu(0) := \min\{m \in \N \mid |f_m(r_0) - a_0 | < 1\}$

          $\nu(0)$ ist wohldefiniert, denn die Menge auf der rechten Seite der Gleicheit ist nichtleer, da nach Voraussetzung $a_0$ ein Häufungspunkt von $(f_n(r_0))_{n \in \N}$ ist und somit eine gegen $r_0$ konvergente Teilfolge existiert.
          Daher besitzt sie aufgrund der natürlichen Wohlordnung der natürlichen Zahlen als nichtleere Teilmenge ein Minimum.

        \item[$b)$] $\nu(n+1) :=  \min\{m \in \N \mid \nu(n) < m \text{ und }  |f_m(r_0) - a_0| < \frac{1}{n+1}\}$

          Wie in $a)$ sieht man ein, dass $\nu(n+1)$ wohldefiniert ist.
      \end{itemize}
      Folglich ist $s_0:= (g_n^0)_{n \in \N} = (f_{\nu(n)})_{n \in \N}$ eine Teilfolge von $(f_n)_{n \in \N}$ und es gilt $\lim_{n \to \infty}(g_n^0(r_0)) = a_0$.

    \item 
      Seien nun $a_n$ und $s_n = (g_m^n)_{m \in \N}$ definiert.
      Es sei nun $a_{n + 1} := \liminf_{m \to \infty} g_m^n(r_{n+1})$.
      Analog zu $1.$ definiere man nun induktiv $s_{n+1} := (g_m^{n + 1})_{m \in \N}$ als Teilfolge von $s_n = (g_m^n)_{m \in \N}$, sodass $\lim_{m \to \infty} g_m^{n + 1}(r_{n+1}) = a_{n + 1}$.
  \end{enumerate}

  Hierauf aufbauend betrachte man nun die Diagonalfolge $s := (g_n^n)_{n \in \N}$.
  Dann ist $s$ eine Teilfolge von $(f_n)_{n \in \N}$ und kofinal zu jeder der Folgen $s_n$.
  Also konvergiert für jedes $n \in \N$ die Folge $s(r_n) = (g_m^m(r_n))_{m \in \N}$ gegen $a_n$.
  Folglich konvergiert für jedes $x \in \Q$ die Folge $s(x) = (g_m^m(x))_{m \in \N}$.

  Es soll nun gezeigt werden, dass $s$ lokal gleichmäßig auf $\R$ konvergiert.
  Sei dazu $x \in \R$ und $\varepsilon > 0$ gegeben.
  Da $G$ gleichgradig stetig ist, existiert eine offene Umgebung $U$ von $x$, sodass
  \begin{displaymath}
    f(U) \subseteq U_\frac{\varepsilon}{3}(f(x))
  \end{displaymath}
  für alle $f \in G$ gilt.
  Da $\Q$ dicht in $\R$ liegt existiert ein $y \in U \cap \Q$.
  Nach Konstruktion konvergiert für alle $g_m^m \in G$ die Folge $g_m^m(y)$.
  Sie ist damit also insbesondere eine \textsc{Cauchy}\hyp{}Folge.
  Es existiert also ein $N \in \N$, sodass für alle $m,n \geq N$
  \begin{displaymath}
    | g_m^m(y) - g_n^n(y) | \leq \frac{\varepsilon}{3}
  \end{displaymath}
  gilt.
  Damit gilt für alle $m, n \geq N$
  \begin{displaymath}
    | g_m^m(x) - g_n^n(x) |
    \leq 
    | g_m^m(x) - g_m^m(y) | +  | g_m^m(y) - g_n^n(y) | +  | g_n^n(y) - g_n^n(x) |
    \leq
    \frac{\varepsilon}{3}.
  \end{displaymath}
  Also ist $s$ eine lokal gleichmäßige \textsc{Cauchy}\hyp{}Folge auf $\R$.
  Da $\R$ vollständig ist, konvergiert $s$ auch lokal gleichmäßig auf $\R$. 
  Aufgrund von Proposition \ref{prop:stetigkonv}(c) konvergiert $s$ auch stetig gegen eine Abbildung $g \colon \R \to \R$ mit $g(r_n) = a_n$ für alle $n \in \N$.
  Daraus folgt $(a)$.
\end{proof}

Zusammenfassend lässt sich also erkennen, dass weder der topologische \textsc{Ascoli}\hyp{}Satz noch der klassische \textsc{Ascoli}\hyp{}Satz in $\mathbf{ZF}$ gelten. Man kann jedoch eine modifizierte Variante des klassischen \textsc{Ascoli}\hyp{}Satzes angeben, welcher sich im Rahmen von $\mathbf{ZF}$ beweisen lässt.

\begin{thm}[Modifizierter \textsc{Ascoli}-Satz]
  Für eine Funktionenmenge $F$ stetiger Abbildungen $f \colon \R \to \R$ sind die folgenden Bedingungen äquivalent:
  \begin{itemize}
    \item[(a)] Jede Folge $(f_n)_{n \in \N}$ in $F$ besitzt eine Teilfolge $(f_{(\nu(n))})_{n \in \N}$ die stetig gegen eine nicht notwendig in $F$ liegende Funktion $g$ konvergiert.

    \item[(b)]
      \begin{itemize}
        \item[($\alpha$)] Für alle $x \in \R$ und jede abzählbare Teilmenge $G \subseteq F$ ist die Menge $G(x) = \{ g(x) \mid g \in G \}$ beschränkt.
        \item[($\beta$)] Jede abzählbare Teilmenge von $F$ ist gleichgradig stetig.
      \end{itemize}
  \end{itemize}
\end{thm}

\begin{proof}
  (a) $\Rightarrow$ (b,$\alpha$):
  Sei $G := \{f_n \in F \mid n \in \N\}$ eine abzählbare Teilmenge von $F$.
  Angenommen $G(x)$ sei unbeschränkt für ein $x \in X$.
  Dann lässt sich für alle $n \in \N$
  \begin{displaymath}
    \nu(n) := \min \{ m \in \N \mid n < |f_m(x)| \}
  \end{displaymath}
  definieren, da zu jeder nichtleeren Teilmenge natürlicher Zahlen ein Minimum existiert.
  Dann ist jedoch die Folge $(f_{\nu(n)})_{n \in \N}$ dergestalt, dass keine Teilfolge von $(f_{\nu(n)}(x))_{n \in \N}$ konvergiert.
  Die Folge $(f_{\nu(n)})_{n \in \N}$ konvergiert somit in $x$ nicht punktweise also nach Proposition \ref{prop:stetigkonv}(1) auch nicht stetig.
  Dies widerspricht jedoch (1).

  (a) $\Rightarrow$ (b,$\beta$):
  Es sei $G$ die der Folge $(f_n)_{n \in \N}$ in $F$ zugrunde liegende Menge.
  Angenommen $G$ für ein $x \in \R$ nicht gleichgradig stetig.
  Dann existiert ein $\varepsilon > 0$, sodass für alle $\delta > 0$ ein $n \in \N$ und $y \in \R$ existiert mit
  \begin{displaymath}
     |x - y| < \delta \quad\text{und}\quad |f_n(x) - f_n(y)| \geq \varepsilon.
  \end{displaymath}
  Für alle $n \in \N$ lassen sich nun
  \begin{align*}
    \nu(n) &:= \min \{ m \in \N \mid \text{ Es gibt ein } y \in [x - 2^{-n}, x + 2^{-n}] \text{ mit } |f_m(x) - f_m(y)| \geq \varepsilon\}, \\
    g_n &:= f_{\nu(n)} \quad\text{und}\\
    x_n &:= \min\{ y \in [x - 2^{-n}, x + 2^{-n}] \mid |f_m(x) - f_m(y)| \geq \varepsilon \}
  \end{align*}
  definieren.
  Dann gelten
  \begin{displaymath}
    \lim_{n \in \N} x_n = x \quad\text{und}\quad|g_n(x) - g_n(x_n)| \geq \varepsilon
  \end{displaymath}
  für alle $n \in \N$.
  Somit kann keine Teilfolge von $(g_n)_{n \in \N}$ stetig in $x$ konvergieren, was jedoch (1) widerspricht.

  (b) $\Rightarrow$ (a):
  Dies folgt bereits aus der entsprechenden Implikation in Satz \ref{thm:classicalascoli}.
\end{proof}

% ...

% Literaturverzeichnis (beginnt auf einer ungeraden Seite)
\cleardoublepage
\nocite{*}
\bibliographystyle{alphadin}
\bibliography{proseminar1516_fabian_gabel}
\addcontentsline{toc}{chapter}{Literaturverzeichnis}

      
% ggf. hier Tabelle mit Symbolen 
% (kann auch auf das Inhaltsverzeichnis folgen)
  
% Stichwortverzeichnis (beginnt auf einer ungeraden Seite)
%\cleardoublepage
%\printindex
%\addcontentsline{toc}{chapter}{Stichwortverzeichnis}
  
% Eidesstattliche Erklaerung.
\cleardoublepage 
% Keine Kopf-/Fußzeile auf dieser Seite
\thispagestyle{empty}

\vspace*{4cm}
\section*{Erklärung}\index{Erklärung}

Hiermit versichere ich, dass ich die vorliegende Arbeit selbständig verfasst habe und alle benutzten Quellen einschließlich der Quellen aus dem Internet und alle sonstigen Hilfsmittel angegeben habe.\vspace{20pt}

\noindent
Ort, den Datum\vspace{30pt}

% Unterschrift (handgeschrieben)

\noindent
Name des Autors

\end{document}
